% Part: sets-functions-relations
% Chapter: relations
% Section: reflections
%
\documentclass[../../../include/open-logic-section]{subfiles}

\begin{document}

\olfileid{sfr}{rel}{ref}
\olsection{తాత్త్విక పరిశీలనలు}

\olref[set]{sec}లో సంబంధాలను కొన్ని సమితులుగా నిర్వచించాం.
ఇక్కడ కాసేపు ఆగి ఒక తాత్త్విక ప్రశ్న అడగాలి: అటువంటి నిర్వచనం
\emph{ఏం చేస్తోంది}? మనం ఏవో అధిభౌతిక తాదాత్మ్య వాస్తవాలను
\emph{కనుగొన్నాం} అని చెప్పాలనుకోవడం చాలా సందేహాస్పదం.
ఉదాహరణకు $\Nat$పై క్రమ సంబంధం అనేది \olref[set]{sec}లో
నిర్వచించిన $R=  \Setabs{\tuple{n,m}}{n, m \in \Nat\text{ మరియు }
n < m}$ అనే సమితియే అని \emph{తేలింది} అనుకోవడం సందేహాస్పదం.
దీనికి మూడు కారణాలు చూద్దాం.

మొదటిది: \olref[set][pai]{wienerkuratowski}లో
$\tuple{a, b} = \{\{a\}, \{a, b\}\}$ అని నిర్వచించాం.
దీనికి బదులు $\lVert a, b\rVert = \{\{b\}, \{a, b\}\} = \tuple{b,a}$
అనే నిర్వచనాన్ని పరిశీలించండి. $a \neq b$ అయినప్పుడు
$\tuple{a, b} \neq \lVert a,b\rVert$.
అయినా క్రమిత జతకు నిర్వచనంగా $\tuple{a,b}$కు బదులు
$\lVert a,b\rVert$ను సమానంగా తీసుకోగలిగేవాళ్లం.
రెండు నిర్వచనాలూ ఒకే విధంగా పనిచేసేవి.
కనుక సహజ సంఖ్యలపై క్రమ సంబంధంగా ``ఉండటానికి'' సమానంగా సరిపోయే
రెండు అభ్యర్థులు ఇప్పుడు ఉన్నాయి:
\begin{align*}
		R &= \Setabs{\tuple{n,m}}{n, m \in \Nat \text{ మరియు }n < m}\\
		S &= \Setabs{\lVert n,m\rVert}{n, m \in \Nat \text{ మరియు }n < m}.
\end{align*}
మూలకాధారిత సమానత్వం ప్రకారం $R \neq S$ కనుక, అవి \emph{రెండూ}
$\Nat$పై క్రమ సంబంధంతో తాదాత్మ్యం కలిగి ఉండలేవు.
కానీ వాస్తవంగా క్రమ సంబంధం $S$ కాక $R$\emph{యే}
(లేదా దానికి విరుద్ధంగా) అని అనడం ఏకపక్షమైన ఎంపిక అవుతుంది;
అందువల్ల కొంత ఇబ్బందికరమూ అవుతుంది.
(ఇది \citeyear{Benacerraf1965}లో బెనసెరాఫ్ ప్రసిద్ధం చేసిన
సమితి-సిద్ధాంత సంక్షేపణవాదానికి వ్యతిరేక వాదనకు చాలా సరళమైన ఉదాహరణ.
దీనిని మరికొన్నిసార్లు తిరిగి పరిశీలిస్తాం.)

రెండవది: \emph{ప్రతి} సంబంధాన్నీ ఒక సమితితో తాదాత్మ్యం చేయాలని భావిస్తే,
సమితి సభ్యత్వ సంబంధమైన $\in$ కూడా ఒక నిర్దిష్ట సమితి కావాలి.
నిజానికి అది $\Setabs{\tuple{x,y}}{x \in y}$ అనే సమితి కావాలి.
కానీ ఈ సమితి ఉందా? రసెల్ వైరుధ్యాన్ని దృష్టిలో ఉంచుకుంటే,
అటువంటి సమితి ఉందనడం సామాన్యమైన వాదన కాదు.
నిజానికి \oliflabeldef{cumul:::part}{\olref[cumul][][]{part}లో
మనం అభివృద్ధి చేసే సమితి సిద్ధాంతం ఈ సమితి ఉనికిని
\emph{నిరాకరిస్తుంది}.\footnote{ముందున్న విషయాన్ని కొద్దిగా ముందే
చెబితే, కారణం ఇది. విరోధ నిరూపణ కోసం
$I = \Setabs{\tuple{x,y}}{x \in y}$ ఉందనుకుందాం.
అప్పుడు $\bigcup \bigcup I$ సార్వత్రిక సమితి అవుతుంది;
ఇది \olref[sfr][z][sep]{thm:NoUniversalSet}కు విరుద్ధం.}}{సమితి
సిద్ధాంతాన్ని కచ్చితమైన స్వీకృతాధార సిద్ధాంతంగా అభివృద్ధి చేయవచ్చు;
ఆ సిద్ధాంతం ఈ సమితి ఉనికిని నిజంగానే నిరాకరిస్తుంది.}
కనుక కొన్ని సంబంధాలను సమితులుగా పరిగణించగలిగినా,
సమితి సభ్యత్వ సంబంధాన్ని ప్రత్యేక సందర్భంగా పరిగణించాల్సి ఉంటుంది.

మూడవది: సంబంధాలను సమితులతో ``తాదాత్మ్యం'' చేసినప్పుడు,
$\tuple{x,y} \in R$కు బదులుగా $Rxy$ అని రాసుకోవచ్చని చెప్పాం.
సభ్యత్వ సంబంధమైన ``$\in$''ను ఒక విధేయంగా \emph{పరిగణిస్తే}
ఇది సమంజసమే. కానీ ``$\in$'' ఒక నిర్దిష్ట రకం సమితిని సూచిస్తుందని
భావిస్తే, ``$\tuple{x,y} \in R$'' అనే పదబంధంలో సమితులను సూచించే
మూడు ఏకవస్తు పదాలు మాత్రమే ఉంటాయి: ``$\tuple{x,y}$'',
``$\in$'', ``$R$''. అటువంటి పేర్ల జాబితి, ``కప్పు కలంపెట్టె బల్ల''
అనే అర్థరహిత సంకేతమాలలాగే, ఒక ప్రవచనాన్ని వ్యక్తపరచలేదు.
మళ్లీ, కొన్ని సంబంధాలను సమితులుగా పరిగణించగలిగినా, సమితి సభ్యత్వ
సంబంధం ప్రత్యేక సందర్భం కావాలి.
(ఇది ఫ్రేగే యొక్క \emph{గుర్రం} భావన వైరుధ్యంలోని ఒక సరళ రూపాన్నీ,
విట్‌గెన్‌స్టైన్ ఒకసారి రసెల్‌కు వ్యతిరేకంగా లేవనెత్తిన ప్రసిద్ధ
అభ్యంతరాన్నీ కలిపి చూపుతోంది.)

మరి దీని నుంచి ఏం తేలుతుంది? సంఖ్యలపై సంబంధాలు సమితులే అని
చెప్పడంలో \emph{తప్పు} ఏమీ లేదు. ఆ మాటను ఏ ఉద్దేశంతో చెబుతున్నామో
అర్థం చేసుకోవాలంతే. మనం ఒక అధిభౌతిక తాదాత్మ్య వాస్తవాన్ని ప్రకటించడం లేదు.
కొన్ని సందర్భాలలో కొన్ని సంబంధాలను కొన్ని సమితులుగా
\emph{పరిగణించవచ్చనీ}, అలా పరిగణిస్తామనీ మాత్రమే గుర్తిస్తున్నాం.

\end{document}
